%% ================================================================
%%  SECTION 3 — Precedent Search
%% ================================================================
\section{Precedent Search}

Everything downstream depends on retrieval. In common law systems, prior rulings are not merely informative; they are binding. Finding the right precedent is the first move a lawyer makes when building an argument, and it is the step from which every other task in the pipeline inherits its quality, or by which it is quietly poisoned. Grounding generation in retrieved case law produces markedly better outputs than letting a language model rely on its parametric memory alone~\cite{nigam-etal-2025-nyayarag}, and pairing retrieval with case-based reasoning pushes legal fidelity further still~\cite{wiratunga-etal-2024-cbrrag}. We take it as fair to say that retrieval is the single highest-leverage point in any legal AI system: get it wrong, and nothing else in the pipeline matters.

\subsection{Keyword-Based and Sparse Retrieval}

For most Indian lawyers, precedent discovery still means typing terms into Indian Kanoon. The platform runs on classical term-matching, and the workhorse behind it is \texttt{BM25}~\cite{rosa-etal-2021-bm25}, a probabilistic ranking function that scores documents by term frequency scaled against inverse document frequency across the corpus, and that we use here as the standard sparse baseline against which neural alternatives are measured. Simple as the formulation sounds, BM25 is stubbornly hard to beat. Empirical comparisons on legal retrieval benchmarks show it matching or outperforming several neural alternatives, particularly when queries contain distinctive legal phrases. We suspect the community often reaches for neural approaches prematurely, without first exhausting what a well-tuned BM25 can deliver. The failure mode, however, is precise and consequential, namely vocabulary mismatch. A query about ``breach of fiduciary duty'' returns nothing when the relevant judgment phrases the same concept as ``violation of trust obligations,'' and no amount of IDF reweighting can bridge a gap where the surface terms simply do not overlap — though, frankly, where exactly the mismatch tips from inconvenient to catastrophic is a line we have not seen the literature pin down with any precision.

\subsection{Dense Retrieval}

Dense retrieval attacks the vocabulary mismatch problem head-on. Queries and documents are encoded into continuous vector representations, and similarity is computed via cosine or dot-product in the embedding space. Transformer-based encoders learn to push semantically related texts close together, so ``breach of fiduciary duty'' and ``violation of trust obligations'' end up in the same neighbourhood despite sharing no surface tokens. Dual-encoder architectures with summarisation-based re-ranking have proven competitive on the COLIEE case law retrieval task~\cite{althammer-etal-2021-dossier}. Scaling this to real corpora, however, is a separate problem. \texttt{FAISS}~\cite{johnson-etal-2021-faiss}, an approximate nearest-neighbour library that supports sub-linear search over billion-scale vector collections on commodity GPUs, supplies the standard infrastructure for dense retrieval at production scale. The cost is non-trivial: index construction and query-time inference both demand significantly more compute than BM25's lightweight inverted indices, and we are honestly not yet sure whether this trade-off pays off for Indian legal corpora, which are large but perhaps not yet at the scale where sparse retrieval visibly collapses.

\subsection{Domain-Adapted Models for Indian Legal Retrieval}

A BERT checkpoint trained on Wikipedia and BooksCorpus knows very little about \emph{ratio decidendi} or ``Section 302 IPC.'' Legal text is unusual enough — long sentences, nested citations, Latin-inflected phrasing — that general-purpose encoders consistently underperform on legal retrieval and classification tasks. The fix is straightforward in principle, namely continued pre-training on in-domain data. \texttt{InLegalBERT}~\cite{paul-etal-2023-inlegalbert}, a BERT-based encoder pre-trained on a large corpus of Supreme Court and High Court judgments, does exactly this and is used as the domain-adapted backbone in most subsequent Indian legal retrieval and classification work. It consistently outperforms both vanilla BERT and multilingual variants across Indian legal benchmarks. We read this result as confirmation that domain shift in legal NLP is not a minor nuisance; it is a first-order problem that general-purpose scale alone does not solve. Whether the right framing here is ``BERT was always going to fail on Indian law'' or ``BERT could have been salvaged with enough adapter training, and we never really tested that alternative at scale,'' we are not, in candour, fully sure.

\subsection{Retrieval-Augmented Generation for Legal Reasoning}

RAG has become the dominant pattern for grounding language model outputs in retrieved evidence, and legal reasoning is a natural fit for it. Instead of relying on whatever a large language model happens to have memorised during pre-training, the generative component receives actual precedents as context. \texttt{NyayaRAG}~\cite{nigam-etal-2025-nyayaanumana}, a retrieval-augmented framework built specifically over Indian Supreme Court and High Court corpora, applies this directly to judgment prediction and shows that retrieval-grounded generation substantially outperforms standalone generation across multiple court levels. \texttt{CBR-RAG}~\cite{wiratunga-etal-2024-cbrrag}, which integrates case-based reasoning with retrieval-augmented generation, goes a step further: rather than retrieving cases that are merely semantically close, it surfaces structurally analogous cases, improving the legal relevance of the context window. These are meaningful advances. However, we note a persistent structural limitation: in every existing Indian legal AI system we have surveyed, retrieved outputs feed into a generation module and stop there. They are never routed onward to structured fact analysis, argument construction, or outcome prediction. Retrieval remains an endpoint rather than an input to a broader reasoning chain — and we have spent some time trying to articulate why this is, without arriving at an answer that fully satisfies us.

%% ================================================================
%%  SECTION 4 — Case Fact Structuring
%% ================================================================
\section{Case Fact Structuring}

An Indian Supreme Court judgment can run to tens of thousands of words. Procedural history, factual narrative, statutory analysis, counsels' submissions, and the final order are all woven together without consistent formatting or explicit section markers. Feed that to a prediction model as raw text, and you are asking it to find a needle in a haystack while blindfolded. We argue that structuring, that is, the decomposition of the document into its functional parts, is a prerequisite rather than a convenience. Three structuring tasks have received sustained attention in Indian legal NLP, namely named entity recognition, rhetorical role classification, and statutory section tagging.

\subsection{Legal Named Entity Recognition}

General-purpose NER (person, location, organisation) is almost useless on a court judgment. The entities that actually matter here are judge names, petitioner and respondent roles, case numbers, statute references, and cited precedents. Kalamkar et al.~\cite{kalamkar-etal-2022-ner} recognised this and built a fourteen-type schema designed specifically for Indian judicial text, covering courts, judges, petitioners, respondents, lawyers, dates, case numbers, statutes, provisions, precedents, witnesses, other persons, organisations, and geographical entities. Transformer-based sequence labelling models were evaluated against this benchmark on Supreme Court and High Court judgments. The downstream payoff is immediate. Precedent entities can be linked to their source judgments for retrieval, party entities feed coreference resolution, and statute entities bridge to the statutory knowledge base. Without this scaffolding, every downstream module is forced to operate over raw token sequences with no structured hooks to latch onto. We find it somewhat surprising that legal NER still receives relatively little attention compared to classification tasks, given how much of the pipeline depends on it — although, on reflection, perhaps the issue is that NER feels solved when in fact it is only solved on the first three or four entity types, and the rest are quietly carrying error that nobody, ourselves included, has audited.

\subsection{Rhetorical Role Classification}

Every sentence in a judgment plays a role: stating facts, citing a prior ruling, framing a legal issue, articulating reasoning, or delivering the operative decision. Rhetorical role labelling assigns these functional categories automatically. \texttt{LegalEval} (SemEval-2023 Task~6)~\cite{modi-etal-2023-semeval} established a multi-court benchmark for this task across Indian jurisdictions, and the results were instructive: automated segmentation is feasible, though far from solved. The hard part is context dependence. ``The appeal is dismissed'' means one thing when it appears inside a recitation of a prior ruling and something entirely different as the operative order of the present case. You cannot classify the sentence in isolation; long-range context is required. Hierarchical and sequential models that propagate information across sentence boundaries help, but we observe that most current approaches still struggle at document lengths typical of Indian high court judgments. Rhetorical segmentation matters because it is essentially the only way to isolate the \emph{ratio decidendi} from the surrounding noise; without it, every sentence carries equal weight, which we keep flagging as the obvious objection but have not yet seen any system in this space fully address.

\subsection{IPC Section Tagging and Statute Linking}

Statute identification asks a deceptively simple question: which sections of the Indian Penal Code (or other legislative instruments) does this judgment actually engage with? \texttt{LegalEval}~\cite{modi-etal-2023-semeval} includes statute identification as one of three subtasks, with annotated data and evaluation protocols. The task matters because it is the bridge between facts and law; without knowing which statutes apply, neither argument construction nor outcome prediction can proceed on solid footing. The persistent difficulty is implicit reference. Judges routinely discuss the elements, interpretation, and applicability of a statute without ever writing down the section number, leaving the model to perform semantic matching between case language and statutory text — which is considerably harder than spotting an explicit ``Section 302'' mention, and where, candidly, the line falls between ``the model has correctly inferred the statute'' and ``the model has guessed plausibly'' is a question we have not yet seen anyone evaluate carefully.

Taken together, NER, rhetorical roles, and statute tagging address complementary dimensions of document structure: who, what role, and under what law. All three should feed into downstream reasoning as intermediate representations. In practice, they do not. Each task has been studied in isolation within Indian legal NLP, and no existing system propagates their combined outputs into a unified pipeline. We are left with a structuring layer that produces terminal outputs, carefully extracted, and then abandoned.

%% ================================================================
%%  SECTION 5 — Argument Graph Construction
%% ================================================================
\section{Argument Graph Construction}

This is the axis where the literature thins out sharply. Legal reasoning is, at its core, argumentative. A decision crystallises from competing claims, supporting evidence, statutory authority, and precedent invoked by analogy or distinction. If you can represent these components, and the relations between them — support, attack, citation, analogy — as a directed graph, then you have made the inferential skeleton of a case explicit and machine-readable. We consider argument graph construction the most structurally ambitious of the five axes in this survey. It is also, by a wide margin, the one with the least Indian legal AI work behind it. What follows therefore draws heavily on the broader argument mining literature and on empirical work conducted in Western jurisdictions.

\subsection{Argument Mining: Foundations}

Argument mining extracts claims, premises, and evidence from text, and maps the relations between them: support, attack, rebuttal. The most comprehensive legal application to date is Habernal et al.'s work on European Court of Human Rights decisions~\cite{habernal-etal-2023-mining}, in which a taxonomy of seventeen argumentative component types was used to annotate and train transformer-based extractors. Performance is reasonable on shorter decisions but degrades noticeably as documents grow longer and argumentative structure becomes more tangled. The field has moved from sentence-level classification (is this sentence a claim or a premise?) to graph-level modelling, where nodes represent argumentative units and directed edges encode support, contradiction, and evidential relations. This shift matters enormously for legal applications. A single precedent might support a petitioner's claim while simultaneously undermining the respondent's, and only a graph can capture that kind of ambivalence — though we should admit we have not actually verified, on Indian data, that ECHR-trained extractors transfer at all, and the longer that omission sits in our own diagnosis, the more it begins to feel like a hole in it.

\subsection{Argument Structure in Legal Corpora}

Almost all empirical work on legal argument mining has been conducted on European and American case law. \texttt{LexGLUE}~\cite{chalkidis-etal-2022-lexglue}, a benchmark suite for legal language understanding that is now the closest thing the field has to a standard, evaluates primarily on ECHR and other Western corpora. Indian courts are simply not represented. That is not a minor gap; it is a blind spot. The structure of a legal argument graph maps naturally onto the components of an Indian judgment: claims correspond to legal issues in dispute, evidence nodes capture trial-established facts, statute nodes provide the normative warrant, and precedent nodes supply analogical support or distinguishing authority. With a complete graph in hand, weak-point detection — the task of identifying claims that lack adequate evidentiary or statutory backing — becomes a matter of graph completeness analysis. We think this capability alone would be worth the annotation investment, even if no further downstream use were envisioned, although we acknowledge we have been making this pitch for some time without anyone funding the annotation, and at some point that pattern itself becomes data we should probably take seriously.

\subsection{The Indian Legal Gap and Graph Database Representations}

As of this writing, no published dataset, annotation schema, or trained model addresses argument graph construction for Indian legal judgments. The gap is total. Several factors explain it. Indian judgments are substantially longer and more digressive than the ECHR decisions that existing argument mining models were designed for. Their argumentative structure is less formulaic; reasoning is often scattered across non-contiguous sections, which makes span extraction unreliable. The most fundamental barrier, however, is the simplest: annotated training data for Indian courts does not exist.

On the representational side, graph databases such as \texttt{Neo4j} are a natural fit: nodes for claims, evidence, statutes, and precedents, with typed edges for support, citation, contradiction, and analogy. The architecture, in some sense, almost designs itself. What makes this axis particularly compelling from a systems perspective is that its inputs already exist, in principle. Retrieved precedents from Section~III, NER entities, rhetorical segments, and statute tags from Section~IV all feed naturally into argument graph construction. This axis is the connective tissue of the pipeline. Without it, the outputs of retrieval and structuring have nowhere coherent to go. We consider Indian legal argument mining a high-priority research direction, and we are frankly surprised it has received so little attention — though, having sat with the surprise long enough, we are starting to think ``surprise'' is no longer the right word, and that perhaps the more concerning state of affairs is that we have stopped being surprised at all.
